\section{Summary of Contributions}
\label{thesis:concl-summary}

The focus of our initial work was on process-based \acp{TEE} such as Intel
\ac{SGX}, ARM TrustZone and Sancus, which allow to protect small portions of
code in order to minimize the \ac{TCB}. In \cref{ch:authexec}, we presented a
framework that allows developers and system administrators to easily develop and
deploy distributed applications on heterogeneous \acp{TEE}. Our framework
leverages an event-driven programming model where each individual component,
called module, can declare input and output endpoints used to exchange messages
between different modules and implement some custom logic in response of an
event. This allows developers to focus exclusively on the application logic,
while the networking, cryptographic operations, \ac{TEE} management and remote
attestation are automatically managed by the framework itself, greatly reducing
the development effort. Another important aspect of our framework is securing
interaction with peripherals thanks to Sancus' secure I/O capability, which
gives strong integrity and authenticity of device drivers and the data exchanged
between a module and an I/O driver. Of course, leveraging \acp{TEE} comes with a
performance overhead; Nevertheless, our evaluation showed that our framework
does not affect user experience in near real-time applications such as smart
home systems.

In recent years, the confidential computing market has been shifting from
process-based to \acs{VM}-based \acp{TEE} such as AMD \sevsnp{} and Intel
\ac{TDX}. This makes development simpler, because applications can run inside
such \acp{TEE} unmodified, but at the same time it greatly increases the
\ac{TCB} and makes management and attestation more challenging. Our contribution
presented in \cref{ch:cloud} discusses these challenges, with a focus on the
boot process of a \ac{CVM} and remote attestation. We showed that booting up a
\ac{CVM} involves several software layers that may be managed by different
entities, and we proposed a classification of incremental \emph{attestation
levels} representing the extent of an attestation in proportion to the whole
boot process. This classification helped us reason about the several steps that
need to be taken to perform remote attestation correctly (and the tools that can
be used to achieve it), as well as the security implications of a \emph{partial}
attestation (i.e., where not all the software components are measured and
attested). Finally, we assessed confidential computing services in public clouds
such as \ac{AWS}, Microsoft Azure and \ac{GCP}, highlighting a crucial
inconsistency between the threat model of \acp{TEE}, where the infrastructure
provider is considered untrusted, and current cloud offerings, where cloud
providers still play an important role in the management and attestation of
\acp{CVM}.

\acp{TEE} isolate code and data of an application from other software in the
same system, thereby protecting their confidentiality and integrity. However,
the application still relies on the untrusted environment (e.g., the host
operating system and the hypervisor) to function properly. Apart from scheduling
operations, applications might need to access system resources such as a clock.
This not only affects availability, which is out of scope of the \ac{TEE} threat
model, but might also lead to confidentiality and integrity breaches depending
on how these resources are actually used. In \cref{ch:time}, we discussed
possible issues related to the use of time inside a \ac{TEE}. First, we provided
a classification of incremental time levels that indicates, for each level, what
guarantees it gives and which attacks are possible. Then, we used this
classification to assess each individual \ac{TEE} according to what support for
trusted time they can provide, and subsequently we identified some common use
cases where time plays a crucial role, such as rate limiting or credential
checking. We highlighted that, for such applications, using an insecure time
source may lead to vulnerabilities that can invalidate the security guarantees
given by \acp{TEE}, and we recommended that developers should take this issue
into account when developing \ac{TEE} applications from scratch or moving
existing workloads to \acp{TEE}.

Finally, \cref{ch:v2x} showcased how \acp{TEE} can be safely integrated into new
protocols and use cases for strong confidentiality, integrity and authenticity
guarantees, while at the same time providing additional benefits such as
increased efficiency and scalability. We demonstrated this in the context of
\ac{ITS} and \ac{V2X}, where previous research had highlighted the need for a
secure and efficient revocation mechanism for vehicle credentials. Learning from
the lessons discussed in \cref{ch:time}, we proposed a novel scheme for
self-revocation that, unlike previous work, does not rely on a trusted time
source in \acp{TEE} but leverages a synchronization protocol between \ac{V2X}
participants based on a logical clock (``epochs''). Compared to related work in
the field, our revocation scheme is the only one that can provide a
\emph{guaranteed} upper bound on revocation time, thanks to our formal
verification work. Furthermore, our evaluation showed that our solution can
handle \ac{V2X} networks with a large number of vehicles and/or revocations.
Indeed, \acp{TEE} allow to partially decentralize credential management, moving
some of the core logic to trusted services running inside vehicles. This allows
novel protocols to be designed, which may provide several benefits compared to
the current state of the art.

In summary, we hope that our research has given a deeper understanding on the
benefits and limitations of \acp{TEE}. Working with this technology is not a
trivial task and presents several technical and operational challenges.
Nevertheless, this solution can provide a robust defense-in-depth technique to
isolate code and data at runtime, reducing the attack surface and the risk of
compromise with relatively low overhead. Crucially, verifiability via remote
attestation is paramount to the security of \ac{TEE} applications to ensure that
isolation mechanisms are correctly in place and the integrity of code and data
has not been compromised. Unfortunately, however, this is not fully achievable
with today's cloud offerings, but we expect that future \ac{TEE} solutions will
make further steps towards real confidential computing.